\documentclass[12pt]{article}
\usepackage[utf8]{inputenc}
\usepackage[T1]{fontenc}
\usepackage{hyperref}
\title{Talysh Language Resources for K-12 Curriculum Development}
\author{Compiled by Claude Sonnet 5 (claude-sonnet-5) for Andr\'as Kornai}
\date{}
\begin{document}
\maketitle

\section{Introduction}
Talysh (ISO 639-3 tly; Glottocode taly1247) is a north-western Iranian language of Azerbaijan and Iran. It is written in a Latin-based alphabet in Azerbaijan and in Perso-Arabic script in Iran. Speaker estimates in the results range from about 230,000 to 800,000. UNESCO is reported to list it as vulnerable. A few textbooks and dictionaries exist, mostly in Russian or Armenian academic publishing.

\section{Primary Educational Resources}
\begin{itemize}
\item Hakob Avchyan and Vardan Voskanyan, \emph{Manual of Talishi} (Yerevan State University Publishing House, 2022, 214 pp.). A textbook with vocabulary and grammar, reviewed in the \emph{Journal of Oriental Studies}. \url{https://journals.ysu.am/j-orient-stud/article/view/11020}
\item Rauf Khalilov, \emph{Learn Talysh}: an English-language beginner's guide, found in a bookseller listing. \url{https://www.betterread.com.au/book/learn-talysh.do}
\end{itemize}
No Iranian-published or Azerbaijani school textbook was found.

\section{Linguistic and Reference Materials}
\begin{itemize}
\item Novruzali Mamedov, \emph{Talyshi-Russian-Azerbaijani Dictionary} (Baku: Nurlan, 2006). The dictionary covers Northern Talysh and has over ten thousand entries. It is described in a review. \url{https://muse.jhu.edu/article/462410}
\item Wikipedia, ``Talysh language''. \url{https://en.wikipedia.org/wiki/Talysh_language}
\item Omniglot, Talysh: scripts and sample text. \url{https://www.omniglot.com/writing/talysh.htm}
\end{itemize}

\section{Religious and Community Resources}
No scripture translations or community organisations were identified in the search results. A Yerevan State University news item discusses Talysh textbook work and assimilation pressure. \url{https://ysu.am/en/news/74871}

\section{Digital and Audio Resources}
No audio, video, corpus or software resources were confirmed.

\section{Curriculum Development Context}
Talysh is not known to be taught in schools. Materials exist for the Northern variety, mainly the Mamedov dictionary and the Yerevan manual. A curriculum would first have to settle which variety and which script to use, given the Latin, Cyrillic and Perso-Arabic traditions.
\end{document}